% ==============================================================
% TSM-A (Nov 13, 2026): Extensions I. Inserted after Results by
% % ==============================================================
% TSM-A (Nov 13, 2026): Extensions I. Inserted after Results by
% % ==============================================================
% TSM-A (Nov 13, 2026): Extensions I. Inserted after Results by
% % ==============================================================
% TSM-A (Nov 13, 2026): Extensions I. Inserted after Results by
% \input{sec_extA} in main.tex. Every number below is read from
% tsm-yield-thresholds/results/tsm_a/ (see figs_tables/numbers.csv);
% scripts in tsm-yield-thresholds/code/tsm_a/.
% ==============================================================
\section{Extensions I: Other Planning Approaches and Other Uncertainty}
\label{sec:extA}

Before a planner trusts the threshold or the cliff check, four questions remain: would a robust, chance-constrained, or risk-averse planner agree; does the check survive uncertain capacity; where does demand put a fab on the map; and how much of the VSS is real when the yield distribution is estimated from a few lots. We answer each with new solves on the base instance (details in Section~\mbox{S-XV}).

\subsection{Robust, Chance-Constrained, and Risk-Averse Plans}
\label{sec:extA-paradigms}

Over the 50 half-ranges of Step~5 we computed (a) a static box-robust plan~\cite{soyster1973}, the deterministic model with every yield low; (b) adjustable (two-stage) robust plans~\cite{bental2004adjustable} that maximize worst-case profit over the eight scenarios, or over the yields only for each demand; (c) budgeted plans~\cite{bertsimas2004} with at most $\Gamma\in\{0,1,2,3\}$ of the yields FE1, FE2, BE low; (d) chance-constrained plans~\cite{charnes1959,luedtke2008sample} that allow an avoidable shortfall with probability at most $\epsilon\in\{0,0.125,0.25\}$; and (e) plans maximizing $(1-\lambda)$ expected profit plus $\lambda$ CVaR$_{0.9}$~\cite{rockafellar2000}. Each plan fixes the wafer starts and is evaluated like the deterministic plan, with the recourse re-optimized in every scenario. Unlike the static box plan, the adjustable robust plans keep the recourse scenario-wise, as the RP does.

\begin{table}[t]
\centering
\caption{Other planning approaches (base instance, 50 half-ranges). Loss: expected profit below the RP (\% of RP). Gain over the deterministic plan in expected profit (\% of RP) or worst-case profit (\% of the deterministic plan's). Gap: unmet demand above the RP's by over one unit.}
\label{tab:extA-paradigms}
\footnotesize
\setlength{\tabcolsep}{2pt}
\begin{tabular}{@{}p{3.05cm}ccccc@{}}
\toprule
 & \textbf{Max} & \textbf{Loss} & \multicolumn{2}{c}{\textbf{Gain $>2\%$ from}} & \textbf{Gap} \\
\cmidrule(lr){4-5}
\textbf{Plan} & \textbf{loss} & \textbf{$>2\%$, $h$} & \textbf{exp.} & \textbf{worst} & \textbf{from} \\
\midrule
Deterministic (EEV) & 8.34 & 0.18 to 0.50 & n/a & n/a & 0.179 \\
RP; box, demand high; robust in yields; $\Gamma=3$; chance$^{a}$; CVaR $\lambda=0.5$ & 0.00 & none & 0.178 & 0.200 & never \\
Box robust, mean demand & 0.19 & none & 0.187 & 0.199 & never \\
Budget $\Gamma=2$ & 0.75 & none & 0.214 & 0.198 & never \\
Budget $\Gamma\in\{0,1\}$ & 4.59 & 0.36 to 0.50 & 0.266 & 0.196 & 0.333 \\
Max-min; CVaR $\lambda=1$ & 2.71 & 0.01 to 0.06 & 0.267 & 0.195 & never \\
\bottomrule
\end{tabular}

\vspace{1pt}
{\footnotesize $^{a}$All three $\epsilon$; infeasible from $h=0.46$ ($\epsilon=0$) and at 0.50 ($\epsilon=0.125$).}
\end{table}

Every plan that protects against the low yield is within 0.005\% of the RP's expected profit at every half-range and passes the 2\% gain at the RP's 0.178 (the box plan at mean demand: within 0.19\%, at 0.187; Table~\ref{tab:extA-paradigms}). At the base costs a full hedge is cheap, so these paradigms agree with the RP. Plans that protect profit against low demand as well behave differently: the max-min plan (equal here to CVaR with $\lambda=1$) starts fewer wafers than the deterministic plan (9.17 versus 10.29 at $h=0.01$), because the worst scenarios have low demand, earns less than it up to $h=0.10$, and gains 2\% only from 0.267; $\Gamma\le1$ plans ($\Gamma=1$ starts the same wafers as $\Gamma=0$, which ignores yield risk) gain 2\% only from 0.266. In worst-case profit, every non-deterministic plan gains 2\% between $h=0.195$ and 0.200. At a tenfold wafer cost the RP's crossing moves to 0.256 and the box plans stay within 0.09\% of the RP. The threshold thus depends on the paradigm, here moving later, not earlier, which supports the third part of our message: the number does not transfer, the hedging mechanism does.

\subsection{Uncertain Capacity}
\label{sec:extA-capacity}

Following the two-stage form of~\cite{geng2009stochastic}, we made the outsourced FE capacity $L_G$, or separately the in-house FE capacity, two-point random ($\pm\delta$, $\delta\le30\%$), independent of yields and demand (16 scenarios). In-house shortfalls cut the wafers processed; the deterministic plan uses mean capacity, split in proportion to fab capacity.

Uncertain outsourcing brings both thresholds in (at $\pm30\%$: 2\% crossing 0.118, gap onset 0.058; all levels in Section~\mbox{S-XV}). Condition~\eqref{eq:cliff} still predicts the onset within 0.00005 once $L_G$ is replaced by its low value; with the mean $L_G$ it would miss by up to 0.102. A fab should compute $\rho$ with the outsourced capacity it can count on, in line with the role of supply uncertainty in outsourcing decisions~\cite{kouvelis2002outsourcing}. Uncertain in-house capacity leaves the deterministic cliff in place (that plan starts 10.50 wafers near the onset, below the reduced 12.60) but limits the stochastic hedge: the RP's own shortfall starts at $h=0.28$ instead of 0.46, and the 2\% crossing moves out to 0.217. The split between fabs now matters: with the solver's split (FE1 first) the deterministic plan loses over 2\% at every spread for $\delta\ge10\%$. Capacity uncertainty thus refines the cliff check rather than breaking it.

\subsection{Demand on the Decision Map}
\label{sec:extA-demand}

Demand forecasts are uncertain, and manufacturers invest in modeling that uncertainty~\cite{wu2010intel}. Fig.~\ref{fig:extA-demand} adds demand to the decision map: demand half-ranges of 0 to 300 units at means $\bar q$ of 675, 750 (base), and 825; the $\bar q=750$ panel reproduces the Step~8 runs.

\begin{figure}[t]
\centering
\includegraphics[width=\columnwidth]{fig_extA_demand_map.pdf}
\caption{\%VSS (color) over FE and demand half-ranges at three demand means. Hatched: service gap above one unit. Dots: solved gap onset; dashed: onset by~\eqref{eq:cliff}; crosses: 2\% crossing ($h=0$: exceeded at every spread).}
\label{fig:extA-demand}
\end{figure}

The gap onset ranges from $h=0.061$ ($\bar q=675$, demand half-range 150 or more) to 0.282 ($\bar q=675$, no demand spread), and the 2\% crossing from every spread to 0.269. Condition~\eqref{eq:cliff} locates the onset in all 15 settings within 0.00004: a low mean with a wide spread sizes the deterministic plan for few dies ($\bar N=806$) while $N_{\max}$ stays at 1{,}050, raising $\rho$ from 0.744 to 0.930. A fab can place itself on the demand axis by computing $\rho$ from its own forecast. Demand moves where the cliff lies, not how it arises, which refines the second part of our message.

\subsection{Scenarios Estimated from Limited Data}
\label{sec:extA-data}

We took the Beta FE yield of Section~\ref{sec:calibration} (mean 0.875, standard deviation 0.15; true VSS \$164.82, 2.60\%) as the truth and drew $n$ lot yields, 200 replications per $n$. The scenario set was built by sample average approximation (SAA)~\cite{kleywegt2002} and by bootstrap~\cite{efron1993bootstrap}; the stochastic plan and the deterministic plan (from the sample mean) were evaluated exactly under the truth.

The optimizer's curse~\cite{smith2006curse} is clear: with five lots the stochastic plan's in-sample profit overstates its realized profit by \$42.1, and it realizes a VSS of \$148.1, 90\% of the attainable (all sample sizes in Section~\mbox{S-XV}). The in-sample VSS is not significantly above the realized VSS at any $n$, since the deterministic plan is overstated on the same sample by almost as much, and in all 2{,}000 replications the stochastic plan beat the deterministic plan under the truth. The larger risk is the decision: with a true \%VSS of 2.60\%, only 38\% of five-lot and 73\% of twenty-lot samples show more than 2\%, because small samples understate the spread of this skewed distribution. Nor can a fab bound its error from its own data: the gap bound of~\cite{mak1999monte,bayraksan2006assessing} covers the true gap in at least 96\% of replications with batches from the truth, but in 26.5\% to 95\% with bootstrap batches. Scenario-based planning is thus only as accurate as its data: Step~2 of Section~\ref{sec:howto} needs a long lot history (here 87\% correct verdicts at 50 lots), while the cliff check of Step~1 needs no scenario set.

Overall, the extensions support the message on the mechanism and the cliff check, refine it on the numbers, and limit it where yield data are scarce.
 in main.tex. Every number below is read from
% tsm-yield-thresholds/results/tsm_a/ (see figs_tables/numbers.csv);
% scripts in tsm-yield-thresholds/code/tsm_a/.
% ==============================================================
\section{Extensions I: Other Planning Approaches and Other Uncertainty}
\label{sec:extA}

Before a planner trusts the threshold or the cliff check, four questions remain: would a robust, chance-constrained, or risk-averse planner agree; does the check survive uncertain capacity; where does demand put a fab on the map; and how much of the VSS is real when the yield distribution is estimated from a few lots. We answer each with new solves on the base instance (details in Section~\mbox{S-XV}).

\subsection{Robust, Chance-Constrained, and Risk-Averse Plans}
\label{sec:extA-paradigms}

Over the 50 half-ranges of Step~5 we computed (a) a static box-robust plan~\cite{soyster1973}, the deterministic model with every yield low; (b) adjustable (two-stage) robust plans~\cite{bental2004adjustable} that maximize worst-case profit over the eight scenarios, or over the yields only for each demand; (c) budgeted plans~\cite{bertsimas2004} with at most $\Gamma\in\{0,1,2,3\}$ of the yields FE1, FE2, BE low; (d) chance-constrained plans~\cite{charnes1959,luedtke2008sample} that allow an avoidable shortfall with probability at most $\epsilon\in\{0,0.125,0.25\}$; and (e) plans maximizing $(1-\lambda)$ expected profit plus $\lambda$ CVaR$_{0.9}$~\cite{rockafellar2000}. Each plan fixes the wafer starts and is evaluated like the deterministic plan, with the recourse re-optimized in every scenario. Unlike the static box plan, the adjustable robust plans keep the recourse scenario-wise, as the RP does.

\begin{table}[t]
\centering
\caption{Other planning approaches (base instance, 50 half-ranges). Loss: expected profit below the RP (\% of RP). Gain over the deterministic plan in expected profit (\% of RP) or worst-case profit (\% of the deterministic plan's). Gap: unmet demand above the RP's by over one unit.}
\label{tab:extA-paradigms}
\footnotesize
\setlength{\tabcolsep}{2pt}
\begin{tabular}{@{}p{3.05cm}ccccc@{}}
\toprule
 & \textbf{Max} & \textbf{Loss} & \multicolumn{2}{c}{\textbf{Gain $>2\%$ from}} & \textbf{Gap} \\
\cmidrule(lr){4-5}
\textbf{Plan} & \textbf{loss} & \textbf{$>2\%$, $h$} & \textbf{exp.} & \textbf{worst} & \textbf{from} \\
\midrule
Deterministic (EEV) & 8.34 & 0.18 to 0.50 & n/a & n/a & 0.179 \\
RP; box, demand high; robust in yields; $\Gamma=3$; chance$^{a}$; CVaR $\lambda=0.5$ & 0.00 & none & 0.178 & 0.200 & never \\
Box robust, mean demand & 0.19 & none & 0.187 & 0.199 & never \\
Budget $\Gamma=2$ & 0.75 & none & 0.214 & 0.198 & never \\
Budget $\Gamma\in\{0,1\}$ & 4.59 & 0.36 to 0.50 & 0.266 & 0.196 & 0.333 \\
Max-min; CVaR $\lambda=1$ & 2.71 & 0.01 to 0.06 & 0.267 & 0.195 & never \\
\bottomrule
\end{tabular}

\vspace{1pt}
{\footnotesize $^{a}$All three $\epsilon$; infeasible from $h=0.46$ ($\epsilon=0$) and at 0.50 ($\epsilon=0.125$).}
\end{table}

Every plan that protects against the low yield is within 0.005\% of the RP's expected profit at every half-range and passes the 2\% gain at the RP's 0.178 (the box plan at mean demand: within 0.19\%, at 0.187; Table~\ref{tab:extA-paradigms}). At the base costs a full hedge is cheap, so these paradigms agree with the RP. Plans that protect profit against low demand as well behave differently: the max-min plan (equal here to CVaR with $\lambda=1$) starts fewer wafers than the deterministic plan (9.17 versus 10.29 at $h=0.01$), because the worst scenarios have low demand, earns less than it up to $h=0.10$, and gains 2\% only from 0.267; $\Gamma\le1$ plans ($\Gamma=1$ starts the same wafers as $\Gamma=0$, which ignores yield risk) gain 2\% only from 0.266. In worst-case profit, every non-deterministic plan gains 2\% between $h=0.195$ and 0.200. At a tenfold wafer cost the RP's crossing moves to 0.256 and the box plans stay within 0.09\% of the RP. The threshold thus depends on the paradigm, here moving later, not earlier, which supports the third part of our message: the number does not transfer, the hedging mechanism does.

\subsection{Uncertain Capacity}
\label{sec:extA-capacity}

Following the two-stage form of~\cite{geng2009stochastic}, we made the outsourced FE capacity $L_G$, or separately the in-house FE capacity, two-point random ($\pm\delta$, $\delta\le30\%$), independent of yields and demand (16 scenarios). In-house shortfalls cut the wafers processed; the deterministic plan uses mean capacity, split in proportion to fab capacity.

Uncertain outsourcing brings both thresholds in (at $\pm30\%$: 2\% crossing 0.118, gap onset 0.058; all levels in Section~\mbox{S-XV}). Condition~\eqref{eq:cliff} still predicts the onset within 0.00005 once $L_G$ is replaced by its low value; with the mean $L_G$ it would miss by up to 0.102. A fab should compute $\rho$ with the outsourced capacity it can count on, in line with the role of supply uncertainty in outsourcing decisions~\cite{kouvelis2002outsourcing}. Uncertain in-house capacity leaves the deterministic cliff in place (that plan starts 10.50 wafers near the onset, below the reduced 12.60) but limits the stochastic hedge: the RP's own shortfall starts at $h=0.28$ instead of 0.46, and the 2\% crossing moves out to 0.217. The split between fabs now matters: with the solver's split (FE1 first) the deterministic plan loses over 2\% at every spread for $\delta\ge10\%$. Capacity uncertainty thus refines the cliff check rather than breaking it.

\subsection{Demand on the Decision Map}
\label{sec:extA-demand}

Demand forecasts are uncertain, and manufacturers invest in modeling that uncertainty~\cite{wu2010intel}. Fig.~\ref{fig:extA-demand} adds demand to the decision map: demand half-ranges of 0 to 300 units at means $\bar q$ of 675, 750 (base), and 825; the $\bar q=750$ panel reproduces the Step~8 runs.

\begin{figure}[t]
\centering
\includegraphics[width=\columnwidth]{fig_extA_demand_map.pdf}
\caption{\%VSS (color) over FE and demand half-ranges at three demand means. Hatched: service gap above one unit. Dots: solved gap onset; dashed: onset by~\eqref{eq:cliff}; crosses: 2\% crossing ($h=0$: exceeded at every spread).}
\label{fig:extA-demand}
\end{figure}

The gap onset ranges from $h=0.061$ ($\bar q=675$, demand half-range 150 or more) to 0.282 ($\bar q=675$, no demand spread), and the 2\% crossing from every spread to 0.269. Condition~\eqref{eq:cliff} locates the onset in all 15 settings within 0.00004: a low mean with a wide spread sizes the deterministic plan for few dies ($\bar N=806$) while $N_{\max}$ stays at 1{,}050, raising $\rho$ from 0.744 to 0.930. A fab can place itself on the demand axis by computing $\rho$ from its own forecast. Demand moves where the cliff lies, not how it arises, which refines the second part of our message.

\subsection{Scenarios Estimated from Limited Data}
\label{sec:extA-data}

We took the Beta FE yield of Section~\ref{sec:calibration} (mean 0.875, standard deviation 0.15; true VSS \$164.82, 2.60\%) as the truth and drew $n$ lot yields, 200 replications per $n$. The scenario set was built by sample average approximation (SAA)~\cite{kleywegt2002} and by bootstrap~\cite{efron1993bootstrap}; the stochastic plan and the deterministic plan (from the sample mean) were evaluated exactly under the truth.

The optimizer's curse~\cite{smith2006curse} is clear: with five lots the stochastic plan's in-sample profit overstates its realized profit by \$42.1, and it realizes a VSS of \$148.1, 90\% of the attainable (all sample sizes in Section~\mbox{S-XV}). The in-sample VSS is not significantly above the realized VSS at any $n$, since the deterministic plan is overstated on the same sample by almost as much, and in all 2{,}000 replications the stochastic plan beat the deterministic plan under the truth. The larger risk is the decision: with a true \%VSS of 2.60\%, only 38\% of five-lot and 73\% of twenty-lot samples show more than 2\%, because small samples understate the spread of this skewed distribution. Nor can a fab bound its error from its own data: the gap bound of~\cite{mak1999monte,bayraksan2006assessing} covers the true gap in at least 96\% of replications with batches from the truth, but in 26.5\% to 95\% with bootstrap batches. Scenario-based planning is thus only as accurate as its data: Step~2 of Section~\ref{sec:howto} needs a long lot history (here 87\% correct verdicts at 50 lots), while the cliff check of Step~1 needs no scenario set.

Overall, the extensions support the message on the mechanism and the cliff check, refine it on the numbers, and limit it where yield data are scarce.
 in main.tex. Every number below is read from
% tsm-yield-thresholds/results/tsm_a/ (see figs_tables/numbers.csv);
% scripts in tsm-yield-thresholds/code/tsm_a/.
% ==============================================================
\section{Extensions I: Other Planning Approaches and Other Uncertainty}
\label{sec:extA}

Before a planner trusts the threshold or the cliff check, four questions remain: would a robust, chance-constrained, or risk-averse planner agree; does the check survive uncertain capacity; where does demand put a fab on the map; and how much of the VSS is real when the yield distribution is estimated from a few lots. We answer each with new solves on the base instance (details in Section~\mbox{S-XV}).

\subsection{Robust, Chance-Constrained, and Risk-Averse Plans}
\label{sec:extA-paradigms}

Over the 50 half-ranges of Step~5 we computed (a) a static box-robust plan~\cite{soyster1973}, the deterministic model with every yield low; (b) adjustable (two-stage) robust plans~\cite{bental2004adjustable} that maximize worst-case profit over the eight scenarios, or over the yields only for each demand; (c) budgeted plans~\cite{bertsimas2004} with at most $\Gamma\in\{0,1,2,3\}$ of the yields FE1, FE2, BE low; (d) chance-constrained plans~\cite{charnes1959,luedtke2008sample} that allow an avoidable shortfall with probability at most $\epsilon\in\{0,0.125,0.25\}$; and (e) plans maximizing $(1-\lambda)$ expected profit plus $\lambda$ CVaR$_{0.9}$~\cite{rockafellar2000}. Each plan fixes the wafer starts and is evaluated like the deterministic plan, with the recourse re-optimized in every scenario. Unlike the static box plan, the adjustable robust plans keep the recourse scenario-wise, as the RP does.

\begin{table}[t]
\centering
\caption{Other planning approaches (base instance, 50 half-ranges). Loss: expected profit below the RP (\% of RP). Gain over the deterministic plan in expected profit (\% of RP) or worst-case profit (\% of the deterministic plan's). Gap: unmet demand above the RP's by over one unit.}
\label{tab:extA-paradigms}
\footnotesize
\setlength{\tabcolsep}{2pt}
\begin{tabular}{@{}p{3.05cm}ccccc@{}}
\toprule
 & \textbf{Max} & \textbf{Loss} & \multicolumn{2}{c}{\textbf{Gain $>2\%$ from}} & \textbf{Gap} \\
\cmidrule(lr){4-5}
\textbf{Plan} & \textbf{loss} & \textbf{$>2\%$, $h$} & \textbf{exp.} & \textbf{worst} & \textbf{from} \\
\midrule
Deterministic (EEV) & 8.34 & 0.18 to 0.50 & n/a & n/a & 0.179 \\
RP; box, demand high; robust in yields; $\Gamma=3$; chance$^{a}$; CVaR $\lambda=0.5$ & 0.00 & none & 0.178 & 0.200 & never \\
Box robust, mean demand & 0.19 & none & 0.187 & 0.199 & never \\
Budget $\Gamma=2$ & 0.75 & none & 0.214 & 0.198 & never \\
Budget $\Gamma\in\{0,1\}$ & 4.59 & 0.36 to 0.50 & 0.266 & 0.196 & 0.333 \\
Max-min; CVaR $\lambda=1$ & 2.71 & 0.01 to 0.06 & 0.267 & 0.195 & never \\
\bottomrule
\end{tabular}

\vspace{1pt}
{\footnotesize $^{a}$All three $\epsilon$; infeasible from $h=0.46$ ($\epsilon=0$) and at 0.50 ($\epsilon=0.125$).}
\end{table}

Every plan that protects against the low yield is within 0.005\% of the RP's expected profit at every half-range and passes the 2\% gain at the RP's 0.178 (the box plan at mean demand: within 0.19\%, at 0.187; Table~\ref{tab:extA-paradigms}). At the base costs a full hedge is cheap, so these paradigms agree with the RP. Plans that protect profit against low demand as well behave differently: the max-min plan (equal here to CVaR with $\lambda=1$) starts fewer wafers than the deterministic plan (9.17 versus 10.29 at $h=0.01$), because the worst scenarios have low demand, earns less than it up to $h=0.10$, and gains 2\% only from 0.267; $\Gamma\le1$ plans ($\Gamma=1$ starts the same wafers as $\Gamma=0$, which ignores yield risk) gain 2\% only from 0.266. In worst-case profit, every non-deterministic plan gains 2\% between $h=0.195$ and 0.200. At a tenfold wafer cost the RP's crossing moves to 0.256 and the box plans stay within 0.09\% of the RP. The threshold thus depends on the paradigm, here moving later, not earlier, which supports the third part of our message: the number does not transfer, the hedging mechanism does.

\subsection{Uncertain Capacity}
\label{sec:extA-capacity}

Following the two-stage form of~\cite{geng2009stochastic}, we made the outsourced FE capacity $L_G$, or separately the in-house FE capacity, two-point random ($\pm\delta$, $\delta\le30\%$), independent of yields and demand (16 scenarios). In-house shortfalls cut the wafers processed; the deterministic plan uses mean capacity, split in proportion to fab capacity.

Uncertain outsourcing brings both thresholds in (at $\pm30\%$: 2\% crossing 0.118, gap onset 0.058; all levels in Section~\mbox{S-XV}). Condition~\eqref{eq:cliff} still predicts the onset within 0.00005 once $L_G$ is replaced by its low value; with the mean $L_G$ it would miss by up to 0.102. A fab should compute $\rho$ with the outsourced capacity it can count on, in line with the role of supply uncertainty in outsourcing decisions~\cite{kouvelis2002outsourcing}. Uncertain in-house capacity leaves the deterministic cliff in place (that plan starts 10.50 wafers near the onset, below the reduced 12.60) but limits the stochastic hedge: the RP's own shortfall starts at $h=0.28$ instead of 0.46, and the 2\% crossing moves out to 0.217. The split between fabs now matters: with the solver's split (FE1 first) the deterministic plan loses over 2\% at every spread for $\delta\ge10\%$. Capacity uncertainty thus refines the cliff check rather than breaking it.

\subsection{Demand on the Decision Map}
\label{sec:extA-demand}

Demand forecasts are uncertain, and manufacturers invest in modeling that uncertainty~\cite{wu2010intel}. Fig.~\ref{fig:extA-demand} adds demand to the decision map: demand half-ranges of 0 to 300 units at means $\bar q$ of 675, 750 (base), and 825; the $\bar q=750$ panel reproduces the Step~8 runs.

\begin{figure}[t]
\centering
\includegraphics[width=\columnwidth]{fig_extA_demand_map.pdf}
\caption{\%VSS (color) over FE and demand half-ranges at three demand means. Hatched: service gap above one unit. Dots: solved gap onset; dashed: onset by~\eqref{eq:cliff}; crosses: 2\% crossing ($h=0$: exceeded at every spread).}
\label{fig:extA-demand}
\end{figure}

The gap onset ranges from $h=0.061$ ($\bar q=675$, demand half-range 150 or more) to 0.282 ($\bar q=675$, no demand spread), and the 2\% crossing from every spread to 0.269. Condition~\eqref{eq:cliff} locates the onset in all 15 settings within 0.00004: a low mean with a wide spread sizes the deterministic plan for few dies ($\bar N=806$) while $N_{\max}$ stays at 1{,}050, raising $\rho$ from 0.744 to 0.930. A fab can place itself on the demand axis by computing $\rho$ from its own forecast. Demand moves where the cliff lies, not how it arises, which refines the second part of our message.

\subsection{Scenarios Estimated from Limited Data}
\label{sec:extA-data}

We took the Beta FE yield of Section~\ref{sec:calibration} (mean 0.875, standard deviation 0.15; true VSS \$164.82, 2.60\%) as the truth and drew $n$ lot yields, 200 replications per $n$. The scenario set was built by sample average approximation (SAA)~\cite{kleywegt2002} and by bootstrap~\cite{efron1993bootstrap}; the stochastic plan and the deterministic plan (from the sample mean) were evaluated exactly under the truth.

The optimizer's curse~\cite{smith2006curse} is clear: with five lots the stochastic plan's in-sample profit overstates its realized profit by \$42.1, and it realizes a VSS of \$148.1, 90\% of the attainable (all sample sizes in Section~\mbox{S-XV}). The in-sample VSS is not significantly above the realized VSS at any $n$, since the deterministic plan is overstated on the same sample by almost as much, and in all 2{,}000 replications the stochastic plan beat the deterministic plan under the truth. The larger risk is the decision: with a true \%VSS of 2.60\%, only 38\% of five-lot and 73\% of twenty-lot samples show more than 2\%, because small samples understate the spread of this skewed distribution. Nor can a fab bound its error from its own data: the gap bound of~\cite{mak1999monte,bayraksan2006assessing} covers the true gap in at least 96\% of replications with batches from the truth, but in 26.5\% to 95\% with bootstrap batches. Scenario-based planning is thus only as accurate as its data: Step~2 of Section~\ref{sec:howto} needs a long lot history (here 87\% correct verdicts at 50 lots), while the cliff check of Step~1 needs no scenario set.

Overall, the extensions support the message on the mechanism and the cliff check, refine it on the numbers, and limit it where yield data are scarce.
 in main.tex. Every number below is read from
% tsm-yield-thresholds/results/tsm_a/ (see figs_tables/numbers.csv);
% scripts in tsm-yield-thresholds/code/tsm_a/.
% ==============================================================
\section{Extensions I: Other Planning Approaches and Other Uncertainty}
\label{sec:extA}

Before a planner trusts the threshold or the cliff check, four questions remain: would a robust, chance-constrained, or risk-averse planner agree; does the check survive uncertain capacity; where does demand put a fab on the map; and how much of the VSS is real when the yield distribution is estimated from a few lots. We answer each with new solves on the base instance (details in Section~\mbox{S-XV}).

\subsection{Robust, Chance-Constrained, and Risk-Averse Plans}
\label{sec:extA-paradigms}

Over the 50 half-ranges of Step~5 we computed (a) a static box-robust plan~\cite{soyster1973}, the deterministic model with every yield low; (b) adjustable (two-stage) robust plans~\cite{bental2004adjustable} that maximize worst-case profit over the eight scenarios, or over the yields only for each demand; (c) budgeted plans~\cite{bertsimas2004} with at most $\Gamma\in\{0,1,2,3\}$ of the yields FE1, FE2, BE low; (d) chance-constrained plans~\cite{charnes1959,luedtke2008sample} that allow an avoidable shortfall with probability at most $\epsilon\in\{0,0.125,0.25\}$; and (e) plans maximizing $(1-\lambda)$ expected profit plus $\lambda$ CVaR$_{0.9}$~\cite{rockafellar2000}. Each plan fixes the wafer starts and is evaluated like the deterministic plan, with the recourse re-optimized in every scenario. Unlike the static box plan, the adjustable robust plans keep the recourse scenario-wise, as the RP does.

\begin{table}[t]
\centering
\caption{Other planning approaches (base instance, 50 half-ranges). Loss: expected profit below the RP (\% of RP). Gain over the deterministic plan in expected profit (\% of RP) or worst-case profit (\% of the deterministic plan's). Gap: unmet demand above the RP's by over one unit.}
\label{tab:extA-paradigms}
\footnotesize
\setlength{\tabcolsep}{2pt}
\begin{tabular}{@{}p{3.05cm}ccccc@{}}
\toprule
 & \textbf{Max} & \textbf{Loss} & \multicolumn{2}{c}{\textbf{Gain $>2\%$ from}} & \textbf{Gap} \\
\cmidrule(lr){4-5}
\textbf{Plan} & \textbf{loss} & \textbf{$>2\%$, $h$} & \textbf{exp.} & \textbf{worst} & \textbf{from} \\
\midrule
Deterministic (EEV) & 8.34 & 0.18 to 0.50 & n/a & n/a & 0.179 \\
RP; box, demand high; robust in yields; $\Gamma=3$; chance$^{a}$; CVaR $\lambda=0.5$ & 0.00 & none & 0.178 & 0.200 & never \\
Box robust, mean demand & 0.19 & none & 0.187 & 0.199 & never \\
Budget $\Gamma=2$ & 0.75 & none & 0.214 & 0.198 & never \\
Budget $\Gamma\in\{0,1\}$ & 4.59 & 0.36 to 0.50 & 0.266 & 0.196 & 0.333 \\
Max-min; CVaR $\lambda=1$ & 2.71 & 0.01 to 0.06 & 0.267 & 0.195 & never \\
\bottomrule
\end{tabular}

\vspace{1pt}
{\footnotesize $^{a}$All three $\epsilon$; infeasible from $h=0.46$ ($\epsilon=0$) and at 0.50 ($\epsilon=0.125$).}
\end{table}

Every plan that protects against the low yield is within 0.005\% of the RP's expected profit at every half-range and passes the 2\% gain at the RP's 0.178 (the box plan at mean demand: within 0.19\%, at 0.187; Table~\ref{tab:extA-paradigms}). At the base costs a full hedge is cheap, so these paradigms agree with the RP. Plans that protect profit against low demand as well behave differently: the max-min plan (equal here to CVaR with $\lambda=1$) starts fewer wafers than the deterministic plan (9.17 versus 10.29 at $h=0.01$), because the worst scenarios have low demand, earns less than it up to $h=0.10$, and gains 2\% only from 0.267; $\Gamma\le1$ plans ($\Gamma=1$ starts the same wafers as $\Gamma=0$, which ignores yield risk) gain 2\% only from 0.266. In worst-case profit, every non-deterministic plan gains 2\% between $h=0.195$ and 0.200. At a tenfold wafer cost the RP's crossing moves to 0.256 and the box plans stay within 0.09\% of the RP. The threshold thus depends on the paradigm, here moving later, not earlier, which supports the third part of our message: the number does not transfer, the hedging mechanism does.

\subsection{Uncertain Capacity}
\label{sec:extA-capacity}

Following the two-stage form of~\cite{geng2009stochastic}, we made the outsourced FE capacity $L_G$, or separately the in-house FE capacity, two-point random ($\pm\delta$, $\delta\le30\%$), independent of yields and demand (16 scenarios). In-house shortfalls cut the wafers processed; the deterministic plan uses mean capacity, split in proportion to fab capacity.

Uncertain outsourcing brings both thresholds in (at $\pm30\%$: 2\% crossing 0.118, gap onset 0.058; all levels in Section~\mbox{S-XV}). Condition~\eqref{eq:cliff} still predicts the onset within 0.00005 once $L_G$ is replaced by its low value; with the mean $L_G$ it would miss by up to 0.102. A fab should compute $\rho$ with the outsourced capacity it can count on, in line with the role of supply uncertainty in outsourcing decisions~\cite{kouvelis2002outsourcing}. Uncertain in-house capacity leaves the deterministic cliff in place (that plan starts 10.50 wafers near the onset, below the reduced 12.60) but limits the stochastic hedge: the RP's own shortfall starts at $h=0.28$ instead of 0.46, and the 2\% crossing moves out to 0.217. The split between fabs now matters: with the solver's split (FE1 first) the deterministic plan loses over 2\% at every spread for $\delta\ge10\%$. Capacity uncertainty thus refines the cliff check rather than breaking it.

\subsection{Demand on the Decision Map}
\label{sec:extA-demand}

Demand forecasts are uncertain, and manufacturers invest in modeling that uncertainty~\cite{wu2010intel}. Fig.~\ref{fig:extA-demand} adds demand to the decision map: demand half-ranges of 0 to 300 units at means $\bar q$ of 675, 750 (base), and 825; the $\bar q=750$ panel reproduces the Step~8 runs.

\begin{figure}[t]
\centering
\includegraphics[width=\columnwidth]{fig_extA_demand_map.pdf}
\caption{\%VSS (color) over FE and demand half-ranges at three demand means. Hatched: service gap above one unit. Dots: solved gap onset; dashed: onset by~\eqref{eq:cliff}; crosses: 2\% crossing ($h=0$: exceeded at every spread).}
\label{fig:extA-demand}
\end{figure}

The gap onset ranges from $h=0.061$ ($\bar q=675$, demand half-range 150 or more) to 0.282 ($\bar q=675$, no demand spread), and the 2\% crossing from every spread to 0.269. Condition~\eqref{eq:cliff} locates the onset in all 15 settings within 0.00004: a low mean with a wide spread sizes the deterministic plan for few dies ($\bar N=806$) while $N_{\max}$ stays at 1{,}050, raising $\rho$ from 0.744 to 0.930. A fab can place itself on the demand axis by computing $\rho$ from its own forecast. Demand moves where the cliff lies, not how it arises, which refines the second part of our message.

\subsection{Scenarios Estimated from Limited Data}
\label{sec:extA-data}

We took the Beta FE yield of Section~\ref{sec:calibration} (mean 0.875, standard deviation 0.15; true VSS \$164.82, 2.60\%) as the truth and drew $n$ lot yields, 200 replications per $n$. The scenario set was built by sample average approximation (SAA)~\cite{kleywegt2002} and by bootstrap~\cite{efron1993bootstrap}; the stochastic plan and the deterministic plan (from the sample mean) were evaluated exactly under the truth.

The optimizer's curse~\cite{smith2006curse} is clear: with five lots the stochastic plan's in-sample profit overstates its realized profit by \$42.1, and it realizes a VSS of \$148.1, 90\% of the attainable (all sample sizes in Section~\mbox{S-XV}). The in-sample VSS is not significantly above the realized VSS at any $n$, since the deterministic plan is overstated on the same sample by almost as much, and in all 2{,}000 replications the stochastic plan beat the deterministic plan under the truth. The larger risk is the decision: with a true \%VSS of 2.60\%, only 38\% of five-lot and 73\% of twenty-lot samples show more than 2\%, because small samples understate the spread of this skewed distribution. Nor can a fab bound its error from its own data: the gap bound of~\cite{mak1999monte,bayraksan2006assessing} covers the true gap in at least 96\% of replications with batches from the truth, but in 26.5\% to 95\% with bootstrap batches. Scenario-based planning is thus only as accurate as its data: Step~2 of Section~\ref{sec:howto} needs a long lot history (here 87\% correct verdicts at 50 lots), while the cliff check of Step~1 needs no scenario set.

Overall, the extensions support the message on the mechanism and the cliff check, refine it on the numbers, and limit it where yield data are scarce.
